\documentclass{situationist-preprint}

\title{The Existence of a Final Theory}
\msid{SCI-20-a}
\msfamily{SCI-20}
\msdiscipline{Philosophy of Science}
\mstype{Principal manuscript}
\mskeywords{final theory; theory of everything; laws of nature; best system; unification; fundamentality; cognitive closure}

\begin{document}
\maketitle

\begin{abstract}
Is there a theory of everything? We argue that there probably is: a final theory exists, though it may exceed human comprehension. A final theory, as we define it, is a finitely axiomatized true theory, stated in terms of the genuine physical magnitudes, from which every law of nature follows. It is a complete set of laws, not a description of a bottom level of objects, and its existence is therefore compatible with the view that every level of objects has a further level beneath it. On the best-system account of laws, a final theory exists if the genuine magnitudes form a fixed stock and some system summarizing the facts about them is best. We argue that both conditions probably hold, on the evidence of the history of unification and of the stability of the stock of magnitudes across scales, and we answer the objections from effective field theory, from the dappled world and from incompleteness. Finally, we argue that nothing guarantees that the final theory is humanly comprehensible, since the simplicity that makes a system best is measured in the vocabulary of nature, not in the cognitive costs of human beings.
\end{abstract}

\section{The problem}

The hope for a final theory is the hope that the explanatory arrows of physics converge: that the laws of the special sciences follow from the laws of physics, that the laws of physics follow from a smaller set of more general laws, and that the regress ends in a set of principles from which everything else follows and which is not itself explained by deeper principles of the same kind. \citet{Weinberg1992} gave the hope its classic statement and argued that the history of physics supports it. Each major advance, from Newton's unification of terrestrial and celestial mechanics to Maxwell's unification of electricity, magnetism and light and the electroweak unification of the twentieth century, reduced the number of independent laws, and the programme of grand unification aimed to continue the process \citep{GeorgiGlashow1974}.

Three positions oppose the hope. The first holds that nature is inexhaustible: at every scale new phenomena and new laws appear, and there is no end to the process. \citet{Dyson1988} defended a version of this view, and the methodology of effective field theory, on which every physical theory is valid only within a range of energies and is to be replaced by another beyond it, is sometimes taken to support it \citep{Weinberg1979}. The second holds that the world is not governed by a unified system of laws at all, but is a patchwork of domains in each of which different models apply: \citet{Cartwright1999} argued that the laws of physics are true only of the highly idealized, shielded situations that experimenters construct. The third holds that a final theory, if it existed, would be a description of a fundamental level of objects, and that there is no such level.

We argue that a final theory probably exists, and that it may exceed human comprehension. The argument depends on separating two ideas that the third position runs together: a final theory is a complete set of laws, not a theory of the smallest things.

\section{Framework}

\begin{definition}\label{d:genuine}
The \emph{genuine magnitudes} are the sparse physical properties and relations: those that make for objective similarity and figure in the laws, such as mass, charge and spatiotemporal separation.
\end{definition}

\begin{definition}\label{d:final}
A theory is \emph{final} if it is (i) true, (ii) stated in the vocabulary of the genuine magnitudes, (iii) finitely axiomatized, and (iv) such that every law of nature is entailed by it.
\end{definition}

\begin{definition}\label{d:comprehensible}
A theory is \emph{humanly comprehensible} if some human being, with any amount of training and with the aid of any tools that leave understanding to the human, could grasp its axioms and see how its central consequences follow from them.
\end{definition}

\begin{assumption}[Humean laws]\label{a:humean}
The laws of nature are the theorems of the best system: the true deductive system, in the vocabulary of the genuine magnitudes, that best balances simplicity and strength.
\end{assumption}

\begin{assumption}[Fixed stock]\label{a:stock}
The genuine magnitudes instantiated at every level of objects belong to a single set, and that set is finite.
\end{assumption}

Assumption~\ref{a:humean} is the best-system account of laws, which MET-12 defends, in the form given by \citet{Lewis1973,Lewis1994}. Definition~\ref{d:genuine} follows the sparse theory of MET-02 in its scientific version, on which the genuine properties are those physics finds at every scale. Assumption~\ref{a:stock} strengthens the fixed-stock principle of MET-10-a, which holds that the genuine properties instantiated at every level belong to one set, by adding that the set is finite. We argue for it in Section~\ref{s:probably}; it is the assumption on which the word \enquote{probably} rests.

\section{Finality without a fundamental level}

\begin{lemma}\label{l:levels}
The existence of a final theory does not require a fundamental level of objects.
\end{lemma}

\begin{proof}
By Definition~\ref{d:final}, a final theory is a set of laws stated in the vocabulary of the genuine magnitudes. Laws so stated quantify over whatever instantiates the magnitudes, at whatever level, and do not mention levels. If, as MET-10 holds, every level of objects has a finer level, and the magnitudes instantiated at every level are drawn from a fixed stock, a single set of laws about those magnitudes can govern every level. Nothing in the definition requires that the quantifiers of the theory range over a lowest level, or that there be objects whose features are grounded in nothing further.
\end{proof}

The model of MET-10-a illustrates Lemma~\ref{l:levels}. In a cube of homogeneous matter divided without end into smaller cubes, the only genuine properties are mass and spatial relations, and the law that the mass of a region is the sum of the masses of its parts holds at every level. The law is complete for the model although the model has no bottom level. Field theories have the same structure: the field equations govern the field in every region, however small, and the regions have no smallest member. A final theory is final because it leaves no law unstated, not because it describes the last things.

\begin{proposition}\label{p:conditional}
Given Assumptions~\ref{a:humean} and~\ref{a:stock}, if there is a best system and it is finitely axiomatized, a final theory exists.
\end{proposition}

\begin{proof}
By Assumption~\ref{a:humean}, the best system is true and is stated in the vocabulary of the genuine magnitudes, and every law is one of its theorems. If it is finitely axiomatized, it satisfies every clause of Definition~\ref{d:final}.
\end{proof}

Proposition~\ref{p:conditional} shifts the question to whether there is a best system with finitely many axioms. On the best-system account, laws exist only if nature is kind enough that some system is best \citep{Lewis1994}; if no system is best, there are, strictly, no laws, and the question of a final theory lapses with them. Assumption~\ref{a:stock} bears on finite axiomatizability. With finitely many magnitudes, a system that states the laws relating them can be finite, and simplicity favors a finite system over an infinite one of equal strength. With infinitely many magnitudes, each requiring its own laws, no finite system could be strong enough to entail all of them.

\section{Why probably}\label{s:probably}

\begin{theorem}\label{t:probably}
A final theory probably exists.
\end{theorem}

\begin{proof}
By Proposition~\ref{p:conditional}, it suffices that there be a best system with finitely many axioms, given Assumption~\ref{a:stock}. Three considerations make this probable. First, the existence of a best system is what the success of physics presupposes and confirms: physics has repeatedly found systems that are simpler and stronger than their predecessors, and the improvements have converged on a small set of dynamical principles rather than dispersing. Second, the stock of genuine magnitudes has been stable across the scales that physics has explored. From molecules to nuclei to quarks, each new level has instantiated mass, charge, spin and spatiotemporal relations, with the addition of a small number of further charges; the number of independent magnitudes has not grown with each new level, and unification has reduced the number of independent laws among them. Third, the reduction of the special sciences to physics, of which chemistry is the clearest case (SCI-11), shows that the laws of higher levels are entailed by those of lower ones together with conditions, so that the laws of physics are the only candidates for axioms. The inference from these considerations is inductive, and it could fail: nature might, at energies not yet explored, reveal magnitudes without end, or the competition between simplicity and strength might have no winner. The conclusion is therefore that a final theory probably exists, not that it certainly does.
\end{proof}

\section{Comprehensibility}

\begin{proposition}\label{p:comprehension}
It is epistemically possible, and not improbable, that the final theory is not humanly comprehensible.
\end{proposition}

\begin{proof}
By Assumption~\ref{a:humean}, what makes a system best is its balance of simplicity and strength in the vocabulary of the genuine magnitudes. That vocabulary is fixed by nature, not by human concepts, and the simplicity of a system is measured by the length of its axioms in it. Nothing in the account ties either the vocabulary or the length to human cognitive capacities. The genuine magnitudes might be ones that human beings can represent only through long chains of mathematical construction, and the shortest axiomatization in their terms might still be too long, or its central derivations too complex, for any human to survey. There are precedents within physics: the ground-state problem for interacting electrons, though fully stated by known laws, is computationally intractable (SCI-11), and the structure of quantum field theories is grasped by most physicists only through approximation schemes. More generally, the cognitive capacities of human beings were shaped by the demands of a particular environment, and \citet{McGinn1989} and \citet{Chomsky1975} have argued that it would be surprising if they were adequate to every problem that nature poses. Since a best system's simplicity is indexed to nature, there is no mechanism that would ensure that the final theory falls within human reach.
\end{proof}

Proposition~\ref{p:comprehension} does not say that the final theory is incomprehensible; it says that its existence does not depend on its comprehensibility, and that the probability that it exceeds human comprehension is not negligible. A final theory that no human could grasp would still exist, and its laws would still be the laws.

\section{Objections and replies}

\begin{objection}[Infinite descent]
If, as MET-10 holds, there is no fundamental level, then every level is governed by laws grounded in a lower one, and there are infinitely many levels of laws. No finite theory can state them.
\end{objection}

\begin{reply}
The objection confuses levels of objects with levels of laws. By Lemma~\ref{l:levels}, laws about the genuine magnitudes govern every level of objects, and the infinitude of levels does not multiply them. MET-10-a itself states that the descent of levels is compatible with a fixed stock of genuine properties, and denies a descent of properties in which every property is constructed from more basic ones without end. What would defeat a final theory is a descent of magnitudes, which Assumption~\ref{a:stock} excludes and which the record of physics does not support.
\end{reply}

\begin{objection}[Effective field theory]
Every successful theory in physics is an effective theory \citep{Weinberg1979}, valid below some energy scale and replaced by another above it. The natural extrapolation is an infinite tower of effective theories with no final member.
\end{objection}

\begin{reply}
That every theory found so far is effective does not show that every theory is. The methodology of effective field theory is neutral: it describes how a theory valid at low energies relates to whatever lies above, and it is consistent with the tower's ending in a theory valid at all scales. Weinberg, who did more than anyone to develop the methodology, was among the strongest advocates of a final theory \citep{Weinberg1992}. Whether the tower ends is the question of Theorem~\ref{t:probably}, and the stability of the stock of magnitudes is evidence that it does.
\end{reply}

\begin{objection}[The dappled world]
The laws of physics are true only of idealized, shielded systems; of the messy world they are false. There is therefore no true unified theory of the world, final or otherwise \citep{Cartwright1999}.
\end{objection}

\begin{reply}
Cartwright's evidence concerns the models by which laws are applied to complex systems, and SCI-10 grants that models are false, representing through partial similarity. The falsity of models does not transfer to the laws. A law stating how the genuine magnitudes evolve is not a claim that any real system behaves as a shielded one does; it is a claim about the contribution each magnitude makes, which is true in unshielded systems as well, where the contributions combine. The difficulty of computing the combined effect, as in chemistry, is computational, not a sign that the laws are false.
\end{reply}

\begin{objection}[Incompleteness]
A final theory strong enough to describe physics would contain arithmetic, and by Gödel's theorem there would be truths it could not prove. So no theory could be final.
\end{objection}

\begin{reply}
Definition~\ref{d:final} requires that the theory entail every law, not that it prove every mathematical truth. Gödel's theorem shows that no consistent, recursively axiomatized theory containing arithmetic proves all arithmetical truths; it does not show that there are laws of nature beyond those the theory entails. Undecidable sentences of the theory's mathematics are not further laws, and their presence is compatible with the theory's leaving no law unstated.
\end{reply}

\section{What the result does not establish}

The argument does not establish what the final theory says, or that any current theory approximates it. It does not establish that the final theory explains its own truth; SCI-18 argues that science cannot explain the existence of the conditions that make it possible, and among these is the holding of the laws themselves. Nor does it establish that human beings will not find the final theory: Proposition~\ref{p:comprehension} records a possibility, not a prediction.

\section{Conclusion}

A final theory is a complete set of laws about the genuine magnitudes, and its existence does not require a bottom level of objects. On the best-system account, it exists if the magnitudes form a finite stock and some system is best, and the history of unification and the stability of the magnitudes across scales make both probable. Since the simplicity that makes a system best is measured in nature's vocabulary, not in human cognitive costs, the final theory may exceed human comprehension.

\begin{thebibliography}{99}\small
\bibitem[Cartwright(1999)]{Cartwright1999} Cartwright, N. (1999). \emph{The Dappled World: A Study of the Boundaries of Science}. Cambridge University Press.
\bibitem[Chomsky(1975)]{Chomsky1975} Chomsky, N. (1975). \emph{Reflections on Language}. Pantheon.
\bibitem[Dyson(1988)]{Dyson1988} Dyson, F. (1988). \emph{Infinite in All Directions}. Harper \& Row.
\bibitem[Georgi and Glashow(1974)]{GeorgiGlashow1974} Georgi, H. and Glashow, S.~L. (1974). Unity of all elementary-particle forces. \emph{Physical Review Letters}, 32(8), 438--441.
\bibitem[Lewis(1973)]{Lewis1973} Lewis, D. (1973). \emph{Counterfactuals}. Blackwell.
\bibitem[Lewis(1994)]{Lewis1994} Lewis, D. (1994). Humean supervenience debugged. \emph{Mind}, 103(412), 473--490.
\bibitem[McGinn(1989)]{McGinn1989} McGinn, C. (1989). Can we solve the mind--body problem? \emph{Mind}, 98(391), 349--366.
\bibitem[Weinberg(1979)]{Weinberg1979} Weinberg, S. (1979). Phenomenological Lagrangians. \emph{Physica A}, 96(1--2), 327--340.
\bibitem[Weinberg(1992)]{Weinberg1992} Weinberg, S. (1992). \emph{Dreams of a Final Theory}. Pantheon.
\end{thebibliography}

\end{document}
